\documentclass[11pt,letterpaper]{article}
\usepackage[T1]{fontenc}
\usepackage[utf8]{inputenc}
\usepackage{lmodern,microtype}
\usepackage[margin=0.79in,headheight=15pt]{geometry}
\usepackage{parskip,enumitem,longtable,tabularx,booktabs,array,xcolor,fancyhdr,hyperref}
\definecolor{navy}{HTML}{17314C}
\definecolor{teal}{HTML}{087F78}
\hypersetup{colorlinks=true,linkcolor=navy,urlcolor=teal}
\setlength{\emergencystretch}{3em}
\setlist[itemize]{leftmargin=*,nosep,topsep=5pt}
\setlist[enumerate]{leftmargin=*,itemsep=3pt}
\renewcommand{\arraystretch}{1.22}
\pagestyle{fancy}\fancyhf{}
\fancyhead[L]{\small Enterprise Data Modeling}\fancyhead[R]{\small Documentation Suite}
\fancyfoot[C]{\thepage}
\newcommand{\srcbook}{\textit{The Data Model Resource Book}, revised edition, Volume 2}
\newcommand{\ResponseLine}{\par\vspace{1.05em}\noindent\rule{\linewidth}{0.4pt}\par}
\newcommand{\ResponseArea}{\ResponseLine\ResponseLine\ResponseLine}
\begin{document}
\begin{center}
{\LARGE\bfseries\color{navy} Normalization and Functional Dependencies}\par
\vspace{0.25em}{\large Supplementary technical reference}\par
\vspace{0.3em}{\small Original engineering study material \textbullet{} 2026 edition}
\end{center}
\vspace{0.8em}
\section{Scope}
Supplementary modern engineering reference. These practices extend the book's industry modeling examples; they are not presented as verbatim material from the 2001 source.


\section{Concepts and design rules}
\begin{tabularx}{\linewidth}{@{}p{0.27\linewidth} p{0.69\linewidth}@{}}
\toprule
\textbf{Topic} & \textbf{Design rule} \\ 
\midrule
Definitions & Functional dependency X determines Y if each admissible relation instance has one Y value for a fixed X; dependencies must be justified semantically, not inferred from a sample. \\ 
1NF & Define atomic domains appropriate to the workload; avoid repeating column groups while permitting deliberate structured types when justified. \\ 
2NF & Eliminate nonkey dependencies on proper subsets of candidate keys in relations with composite candidate keys. \\ 
3NF & For every nontrivial FD X to A, require X to be a superkey or A to be prime; document the precise definition used. \\ 
BCNF & For every nontrivial functional dependency X to Y, require X to be a superkey; weigh decomposition against dependency preservation. \\ 
Further forms & Analyze multivalued dependencies and join dependencies when they are evidenced; do not normalize indefinitely without a query and correctness objective. \\ 
\bottomrule
\end{tabularx}

\section{Worked method}
\begin{enumerate}
\item Capture a source requirement and distinguish data guarantees from presentation preferences.
\item State formal or operational invariants, define test cases and document exceptions.
\item Identify the simplest implementation that satisfies current correctness and performance needs.
\item Design one adverse case involving concurrent writes, a missing parent or a late-arriving change.
\item Record migration and rollback implications and obtain review sign-off.
\end{enumerate}

\section{Evidence and acceptance}
Evidence should include a glossary, candidate key/FD list, test input-output pairs, relevant query plan where applicable, change history, and model reviewer approval. Negative examples must be recorded along with positive examples.


\vfill
\hrule
{\footnotesize\textbf{Provenance.} Supplementary engineering practice: original content; the cited book provides domain-modeling context.}
\end{document}
